%%% Tento soubor obsahuje definice různých užitečných maker a prostředí %%%
%%% Další makra připisujte sem, ať nepřekáží v ostatních souborech.     %%%
%%% This file contains definitions of various useful macros and environments      %%%
%%% Assign additional macros here so that they do not interfere with other files. %%%

\usepackage{ifpdf}
\usepackage{ifxetex}
\usepackage{ifluatex}

%%% Nastavení pro použití samostatné bibliografické databáze.
%%% Settings for using a separate bibliographic database.
\input biblatex-setup

%% Přepneme na českou sazbu, fonty Latin Modern a kódování češtiny
\ifthenelse{\boolean{xetex}\OR\boolean{luatex}}
   { % use fontspec and OpenType fonts with utf8 engines
			\usepackage[english,slovak,czech]{babel}
			\usepackage[autostyle,english=british,czech=quotes]{csquotes}
			\usepackage{fontspec}
			\defaultfontfeatures{Ligatures=TeX,Scale=MatchLowercase}
   }
   {
			\usepackage[english,slovak,czech]{babel}
			\usepackage{lmodern}
			\usepackage[T1]{fontenc}
			\usepackage{textcomp}
			\usepackage[utf8]{inputenc}
			\usepackage[autostyle,english=british,czech=quotes]{csquotes}
	 }
\ifluatex
\makeatletter
\let\pdfstrcmp\pdf@strcmp
\makeatother
\fi

\usepackage[a-2u]{pdfx}     % výsledné PDF bude ve standardu PDF/A-2u
                            % resulting PDF will be in the PDF / A-2u standard

%%% Další užitečné balíčky (jsou součástí běžných distribucí LaTeXu)
\usepackage{amsmath}        % rozšíření pro sazbu matematiky / extension for math typesetting
\usepackage{amsfonts}       % matematické fonty / mathematical fonts
\usepackage{amssymb}        % symboly / symbols
\usepackage{amsthm}         % sazba vět, definic apod. / typesetting of sentences, definitions, etc.
\usepackage{bm}             % tučné symboly (příkaz \bm) / bold symbols (\bm command)
\usepackage{graphicx}       % vkládání obrázků / graphics inserting
\usepackage{listings}       % vylepšené prostředí pro strojové písmo / improved environment for source codes typesetting
\usepackage{fancyhdr}       % prostředí pohodlnější nastavení hlavy a paty stránek / environment for more comfortable adjustment of the head and foot of the pages
\usepackage{icomma}         % inteligetní čárka v matematickém módu / intelligent comma in math mode
\usepackage{dcolumn}        % lepší zarovnání sloupců v tabulkách / better alignment of columns in tables
\usepackage{booktabs}       % lepší vodorovné linky v tabulkách / better horizontal lines in tables
\usepackage{ltablex}       % tabularx + longtable (tables stay in place and may break across pages)
\keepXColumns
\makeatletter
\@ifpackageloaded{xcolor}{
   \@ifpackagewith{xcolor}{usenames}{}{\PassOptionsToPackage{usenames}{xcolor}}
  }{\usepackage[usenames]{xcolor}} % barevná sazba / color typesetting
\makeatother
\usepackage{multicol}       % práce s více sloupci na stránce / work with multiple columns on a page
\usepackage{multirow}
\usepackage{caption}
\usepackage{sepfootnotes}
\newfootnotes*{tabfn}
% Poznámky pod čarou (sepfootnotes) – definice obsahu
% Klíče jsou stabilní identifikátory, které se používají v textu přes \fn{<klic>}.

\sepfootnotecontent{vymezeni-pm25}{Particulate Matter $2.5\,\mu\text{m}$}
\sepfootnotecontent{vymezeni-pm10}{Particulate Matter $10\,\mu\text{m}$}
\sepfootnotecontent{vymezeni-epa}{www.epa.gov}
\sepfootnotecontent{vymezeni-eco2}{odhadovaná koncentrace CO\textsubscript{2} senzorem\parencite{calvoScalableIoTArchitecture2022}}

% Tabulkové poznámky (sepfootnotes aparát tabfn) – definice obsahu

\tabfnnotecontent{vymezeni-aot40}{\textit{Pro účely tohoto zákona AOT40 znamená součet rozdílů mezi hodinovou koncentrací větší než 80 µg.m\textsuperscript{-3} (= 40 ppb) a hodnotou 80 µg.m\textsuperscript{-3} v dané periodě užitím pouze hodinových hodnot změřených každý den mezi 08:00 a 20:00 SEČ, vypočtený z hodinových hodnot v letním období (1. května - 31. července).} \parencite[příl.~1, pozn.~5]{wolterskluwerZakon20120122012}}

\tabfnnotecontent{vymezeni-la-a-index}{Index „A“ v označení $L_A$ znamená, že jde o hladinu akustického tlaku měřenou s frekvenční váhovou charakteristikou A (A-vážená hladina); např. $LA_{eq,T}$ je A-vážená ekvivalentní hladina akustického tlaku za dobu $T$. \parencite{info@wolterskluwer.czNarizeni2722011}}

\tabfnnotecontent{vymezeni-la-eu-note}{\textit{Poznámka k EU (proč tu nejsou hodnoty):} Směrnice \parencite{SmerniceEvropskehoParlamentu2002} nestanovuje jednotné číselné hygienické limity (ty si určují členské státy). Zavádí ale společné ukazatele hluku: $L_{den}$ (\textit{day-evening-night level}) je vážená průměrná hladina za 24 h, kde se večerní část penalizuje +5 dB a noční část +10 dB; $L_{night}$ (\textit{night level}) je průměrná hladina za noční období používaná hlavně pro hodnocení rušení spánku.
Zdroj: \parencite{info@wolterskluwer.czNarizeni2722011}.}

\newcommand{\fn}[1]{\sepfootnote{#1}}

% Table footnote marks/texts that can be placed apart while preserving numbering.
\makeatletter
\newcommand{\tfnmark}[1]{%
  % tabularx typesets its contents in a "trial" box to measure column widths.
  % In that pass it disables \@footnotetext but still executes \footnotemark,
  % which would step counters / create anchors twice and break numbering/links.
  \ifx\@footnotetext\TX@trial@ftn
    \textsuperscript{0}% dummy mark (trial pass only; not printed in the final table)
  \else
    \@ifundefined{tfn@#1}{%
      \footnotemark%
      \expandafter\xdef\csname tfn@#1\endcsname{\number\value{footnote}}%
    }{%
      \footnotemark[\csname tfn@#1\endcsname]%
    }%
  \fi
}
\newcommand{\tfntext}[1]{%
  \begingroup
    \@ifundefined{tfn@#1}{%
      \PackageWarning{sepfootnotes}{Missing table footnote mark for key `#1'}%
    }{%
      % Create the *destination* for hyperlinks using the same identifier that
      % hyperref expects (Hfootnote.<n>), then typeset the footnote text with the
      % stored printed number (<n>).
      \leavevmode\Hy@raisedlink{\hyper@anchorstart{Hfootnote.\csname tfn@#1\endcsname}\hyper@anchorend}%
      \footnotetext[\csname tfn@#1\endcsname]{\printtabfnnote{#1}}%
    }%
  \endgroup
}
\makeatother

\usepackage{enumitem}
\setlist[itemize]{noitemsep, topsep=0pt, partopsep=0pt}
\setlist[enumerate]{noitemsep, topsep=0pt, partopsep=0pt}
\setlist[description]{noitemsep, topsep=0pt, partopsep=0pt}

\usepackage{tocloft}
\setlength\cftparskip{0pt}
\setlength\cftbeforechapskip{1.5ex}
\setlength\cftfigindent{0pt}
\setlength\cfttabindent{0pt}
\setlength\cftbeforeloftitleskip{0pt}
\setlength\cftbeforelottitleskip{0pt}
\setlength\cftbeforetoctitleskip{0pt}
\renewcommand{\cftlottitlefont}{\Huge\bfseries\sffamily}
\renewcommand{\cftloftitlefont}{\Huge\bfseries\sffamily}
\renewcommand{\cfttoctitlefont}{\Huge\bfseries\sffamily}

% vyznaceni odstavcu
% differentiation of new paragraphs
\parindent=0pt
\parskip=11pt

% zakaz vdov a sirotku - jednoradkovych pocatku ci koncu odstavcu na prechodu mezi strankami
% Prohibition of widows and orphans - single-line beginning and end of paragraph at the transition between pages
\clubpenalty=1000
\widowpenalty=1000
\displaywidowpenalty=1000

% nastaveni radkovani
% setting of line spacing
\renewcommand{\baselinestretch}{1.20}

% nastaveni pro nadpisy - tucne a bezpatkove
% settings for headings - bold and sans serif
\usepackage{sectsty}    
\allsectionsfont{\sffamily}

% nastavení hlavy a paty stránek
% page head and foot settings
\makeatletter
\if@twoside%
    \fancypagestyle{fancyx}{%
			\fancyhf{}                                     
      \fancyhead[RE]{\rightmark}                  
      \fancyhead[LO]{\leftmark}                  
      \fancyfoot[RO,LE]{\thepage}                    
      \renewcommand{\headrulewidth}{.5pt}            
      \renewcommand{\footrulewidth}{.5pt}            
    }
    \fancypagestyle{plain}{%
			\fancyhf{}                                     
    	\fancyfoot[RO,LE]{\thepage}                    
    	\renewcommand{\headrulewidth}{0pt}             
    	\renewcommand{\footrulewidth}{0.5pt}
    }          
\else
    \fancypagestyle{fancyx}{%
			\fancyhf{}                                     
      \fancyhead[R]{\leftmark}                  
      \fancyfoot[R]{\thepage}                    
      \renewcommand{\headrulewidth}{.5pt}            
      \renewcommand{\footrulewidth}{.5pt}            
    }
    \fancypagestyle{plain}{%                       
    	\fancyhf{} % clear all header and footer fields
    	\fancyfoot[R]{\thepage}                    
    	\renewcommand{\headrulewidth}{0pt}             
    	\renewcommand{\footrulewidth}{0.5pt}
    }          
\fi
\renewcommand*{\cleardoublepage}{\clearpage\if@twoside \ifodd\c@page\else
	\hbox{}%
	\thispagestyle{empty}%
	\newpage%
	\if@twocolumn\hbox{}\newpage\fi\fi\fi
}
\makeatother

% Tato makra přesvědčují mírně ošklivým trikem LaTeX, aby hlavičky kapitol
% sázel příčetněji a nevynechával nad nimi spoustu místa. Směle ignorujte.
% These macros convince with a slightly ugly LaTeX trick to make chapter headers
% bet more sane and didn't miss a lot of space above them. Be boldly ignore it.
\makeatletter
\def\@makechapterhead#1{
  {\parindent \z@ \raggedright \sffamily
   \Huge\bfseries \thechapter. #1
   \par\nobreak
   \vskip 20\p@
}}
\def\@makeschapterhead#1{
  {\parindent \z@ \raggedright \sffamily
   \Huge\bfseries #1
   \par\nobreak
   \vskip 20\p@
}}
\makeatother

% Trochu volnější nastavení dělení slov, než je default.
% Slightly looser hyphenation setting than default.
\lefthyphenmin=2
\righthyphenmin=2

% Zapne černé "slimáky" na koncích řádků, které přetekly, abychom si jich lépe všimli.
% Turns on the black "snails" at the ends of the lines that overflowed to get us noticed them better.
\overfullrule=1mm

\def\BibTeX{{\rm B\kern-.05em{\sc i\kern-.025em b}\kern-.08em
    T\kern-.1667em\lower.7ex\hbox{E}\kern-.125emX}}

%% Balíček hyperref, kterým jdou vyrábět klikací odkazy v PDF,
%% ale hlavně ho používáme k uložení metadat do PDF (včetně obsahu).
%% Většinu nastavítek přednastaví balíček pdfx.
%% A hyperref package that can be used to produce clickable links in PDF,
%% but we mainly use it to store metadata in PDF (including content).
%% Most settings are preset by the pdfx package.
\hypersetup{unicode}
\hypersetup{breaklinks=true}
\hypersetup{hidelinks}
\hypersetup{colorlinks=true,urlcolor=blue}

\renewcommand{\UrlBreaks}{\do\/\do\=\do\+\do\-\do\_\do\ \do\a\do\b\do\c\do\d%
\do\e\do\f\do\g\do\h\do\i\do\j\do\k\do\l\do\m\do\n\do\o\do\p\do\q\do\r\do\s%
\do\t\do\u\do\v\do\w\do\x\do\y\do\z\do\A\do\B\do\C\do\D\do\E\do\F\do\G\do\H%
\do\I\do\J\do\K\do\L\do\M\do\N\do\O\do\P\do\Q\do\R\do\S\do\T\do\U\do\V\do\W%
\do\X\do\Y\do\Z\do\1\do\2\do\3\do\4\do\5\do\6\do\7\do\8\do\9\do\0}
\urlstyle{tt}

%%% Prostředí pro sazbu kódu, případně vstupu/výstupu počítačových
%%% programů. (Vyžaduje balíček listings -- fancy verbatim.)
%%% Environment for source code typesetting, or computer input/output
%%% programs. (Requires package listings - fancy verbatim.)
\lstnewenvironment{code}[3]{\lstset{language=#1, caption=#2, label=#3, basicstyle=\small\ttfamily, numbers=left, xleftmargin=2em, frame=single, framexleftmargin=2em}}{}

%%% Tato část obsahuje texty závislé na typu práce, jazyku a pohlaví %%%
%%% This part contains texts depending on the type of work, language and gender %%%

\newcommand{\ifstringequal}[4]{%
  \ifnum\pdfstrcmp{#1}{#2}=0
  #3%
  \else
  #4%
  \fi
}

\def\TypPraceBP{BAKALÁŘSKÁ PRÁCE}
\def\TypPraceDP{DIPLOMOVÁ PRÁCE}
\def\SeznamZkratek{Seznam použitých zkratek}
\def\Prilohy{Přílohy}
\def\VSE{Vysoká škola ekonomická v Praze}
\def\FIS{Fakulta informatiky a statistiky}
\def\StudijniProgramText{Studijní program}
\def\SpecializaceText{Specializace}
\def\AutorText{Autor}
\def\VedouciText{Vedoucí práce}
\def\KonzultantText{Konzultant práce}
\def\Praha{Praha}
\def\PodekovaniText{Poděkování}
\def\ProhlaseniAIText{Prohlášení o využití umělé inteligence}
\def\bibnamex{Použitá literatura}
\renewcommand*{\lstlistlistingname}{Seznam zdrojových kódů}
\def\lstlistingname{Výpis}

\makeatletter
\ifstringequal{\Jazyk}{eng}{\main@language{english}}{}
\ifstringequal{\Jazyk}{slo}{\main@language{slovak}}{}
\makeatother

\ifstringequal{\Jazyk}{eng}{
         \def\TypPraceBP{BACHELOR THESIS}
         \def\TypPraceDP{MASTER THESIS}
				 \def\SeznamZkratek{List of abbreviations}
				 \def\Prilohy{Appendices}
				 \def\VSE{Prague University of Economics and Business}
				 \def\FIS{Faculty of Informatics and Statistics}
				 \def\StudijniProgramText{Study program}
				 \def\SpecializaceText{Specialization}
				 \def\AutorText{Author}
				 \def\VedouciText{Supervisor}
         \def\KonzultantText{Consultant}
				 \def\Praha{Prague}
				 \ifstringequal{\TypPrace}{BP}{\def\Praci{bachelor thesis }}{}
         \ifstringequal{\TypPrace}{DP}{\def\Praci{master thesis }}{}
				 \def\PodekovaniText{Acknowledgements}
				 \def\bibnamex{References}
         \renewcommand*{\lstlistlistingname}{List of source codes}
         \def\lstlistingname{Source code}
				 \ifdef{\finalnamedelim}{\renewcommand*{\finalnamedelim}{\addspace and \addspace}}%
}{}

\ifstringequal{\Jazyk}{slo}{
         \def\TypPraceBP{BAKALÁRSKA PRÁCA}
         \def\TypPraceDP{DIPLOMOVÁ PRÁCA}
				 \def\SeznamZkratek{Zoznam použitých skratiek}
				 \def\Prilohy{Prílohy}
				 \def\StudijniProgramText{Študijný program}
				 \def\SpecializaceText{Špecializácia}
				 \def\VedouciText{Vedúci práce}
         \def\KonzultantText{Konzultant práce}
				 \ifstringequal{\TypPrace}{BP}{\def\Praci{bakalársku prácu }}{}
         \ifstringequal{\TypPrace}{DP}{\def\Praci{diplomovú prácu }}{}
				 \def\PodekovaniText{Poďakovanie}
				 \def\bibnamex{Použitá literatúra}
				 \renewcommand*{\lstlistlistingname}{Zoznam zdrojových kódov}
				 \def\lstlistingname{Výpis}
}{}

\ifstringequal{\TypPrace}{BP}{\def\TypPraceText{\TypPraceBP}}{}
\ifstringequal{\TypPrace}{DP}{\def\TypPraceText{\TypPraceDP}}{}
